% Generated draft: requires mathematical and source review.
\documentclass[11pt]{article}
\usepackage[margin=25mm]{geometry}
\usepackage{fontspec}
\setmainfont{Latin Modern Roman}
\usepackage{amsmath,amssymb,mathtools,unicode-math}
\setmathfont{Latin Modern Math}
\usepackage{enumitem,needspace,xurl,hyperref}
\setlength{\parindent}{0pt}
\setlength{\parskip}{5pt}
\setlength{\emergencystretch}{3em}
\newcommand{\N}{\mathbb{N}}
\newcommand{\Z}{\mathbb{Z}}
\newcommand{\Q}{\mathbb{Q}}
\newcommand{\R}{\mathbb{R}}
\newcommand{\C}{\mathbb{C}}
\newcommand{\F}{\mathbb{F}}
\newcommand{\A}{\mathbb{A}}
\newcommand{\PP}{\mathbb{P}}
\newcommand{\E}{\mathbb{E}}
\newcommand{\eps}{\varepsilon}
\newcommand{\dd}{\,\mathrm{d}}
\begin{document}
\section*{UnsolvedMath ID: 30002712}
\textbf{Source problem code:} catalogue entry
\subsection{Definitions and mathematical statement}
The supplied question is to classify local Oort groups, described as groups for which every local \(G\)-extension lifts to characteristic zero. The source record does not specify the characteristic, the residue field, or the admissible class of groups. The following gives a parameterized formulation under explicitly stated conventions; these conventions require confirmation against the cited sources.

Let \(p\) be a prime, let \(k\) be an algebraically closed field of characteristic \(p\), and let \(G\) be a finite group. In this formulation, a local \(G\)-action is a faithful continuous action
\[
\rho\colon G\hookrightarrow \operatorname{Aut}_{k}(k[[z]]),
\]
where continuity refers to the \(z\)-adic topology. Equivalently, it describes a local Galois extension \(k[[z]]/k[[z]]^{G}\) with specified Galois group \(G\).

Say that \(\rho\) lifts to characteristic zero if there exist a complete discrete valuation ring \(R\) of characteristic zero, with maximal ideal \(\mathfrak m_R\) and a specified identification \(R/\mathfrak m_R\cong k\), a faithful continuous action
\[
\widetilde\rho\colon G\hookrightarrow \operatorname{Aut}_{R}(R[[Z]]),
\]
and a \(G\)-equivariant \(k\)-algebra isomorphism
\[
R[[Z]]/\mathfrak m_R R[[Z]]\cong k[[z]].
\]
Here continuity upstairs refers to the \((\mathfrak m_R,Z)\)-adic topology, and the action on the left is induced by \(\widetilde\rho\). The ring \(R\) and the lifted action may depend on \(\rho\).

For a specified class \(\mathcal C\) of finite groups, define
\[
\mathcal L(k;\mathcal C)=
\left\{G\in\mathcal C:\text{every faithful continuous action of }G\text{ on }k[[z]]\text{ lifts as above}\right\}.
\]
\textbf{Classification problem, subject to the unresolved conventions:} determine \(\mathcal L(k;\mathcal C)\) up to group isomorphism for the intended choices of \(p\), \(k\), and \(\mathcal C\).

The record does not determine whether the lifting property is required over one fixed \(k\), over every algebraically closed field of a fixed positive characteristic, or across all positive characteristics. It also does not specify whether admissible groups must occur as local Galois groups. If a group admits no such local action, the displayed universal lifting property holds vacuously unless an occurrence condition is imposed. These distinctions remain unresolved and prevent identifying a unique classification target from the supplied wording alone.
\subsection{Short English statement}
Classify the finite groups for which every action on a formal disc in positive characteristic can be lifted to an action over a complete discrete valuation ring of characteristic zero. The intended characteristic, residue-field quantifiers, and admissible groups need confirmation. The supplied record labels the classification open; this draft does not independently verify that status.
\subsection{Sources}
\begin{itemize}
\item Supplied local UnsolvedMath registry record, permanent ID \(30002712\), source problem code \texttt{catalogue entry}, title \emph{Classification of Local Oort Groups}, in the Oberwolfach Reports Open Problems collection. The record gives proposed year \(2014\) and names no individual proposer.
\item Supplied report citation: \emph{Valuation Theory and Its Applications} (2015), Oberwolfach Reports, DOI \texttt{10.4171/owr/2014/49}, pp.\ 2757--2823. \url{https://doi.org/10.4171/owr/2014/49}. The supplied background describes these reports as documenting mathematical workshops and their open-problem discussions.
\item Supplied literature reference: \emph{Oberwolfach Report 49/2014}. \url{https://publications.mfo.de/bitstream/handle/mfo/3439/OWR_2014_49.pdf?isAllowed=y&sequence=1}.
\item Supplied literature reference: Chinburg--Guralnick--Harbater, \emph{Global Oort groups}. \url{https://www2.math.upenn.edu/~harbater/oortpaper3.pdf}.
\item Literature assessment dated exactly \texttt{2026-08-21}: the supplied curation labels the problem open and states that cyclic groups are local Oort groups, that several additional cases are known, and that a complete classification remains unresolved. This is supplied dated curation derived from the OWR material, not an independently verified assessment; its status and cited claims require independent verification.
\end{itemize}
\end{document}
